\subsection{特征值的变分刻画}

在本节中，我们假设 K = C。

设 \(u\) 是 E 的紧正自同态，\((\lambda_n(u))_{n\in \mathrm{I_E}}\) 是 \(u\) 的特征值递减序列。我们置 \(\mathrm{I} = \mathrm{I_E}\) 。

对 E 的每个闭子空间 F，我们记

\[
r _ {\mathrm{F}} (u) = \inf _ {x \in \mathrm{F} - \{0 \}} \frac {\langle x \mid u (x) \rangle}{\| x \| ^ {2}}, \quad \mathrm{R} _ {\mathrm{F}} (u) = \sup _ {x \in \mathrm{F} ^ {\circ} - \{0 \}} \frac {\langle x \mid u (x) \rangle}{\| x \| ^ {2}},
\]

其中下确界（相应地，上确界）是在 \([0, +\infty]\) 中取的。

对每个 \(n \in \mathbf{N}\) ，用 \(\mathcal{F}_n\) 表示 E 的向量子空间 \(\mathrm{F} \subset \mathrm{E}\) （维数为 \(n\) ）所成的集。我们称子空间 \(\mathrm{F} \in \mathcal{F}_n\) 适应于 \(u\) ，如果它有一个标准正交基 \((f_i)_{0 \leqslant i \leqslant n-1}\) ，使得 \(u(f_i) = \lambda_i(u)f_i\) （对 \(0 \leqslant i \leqslant n-1\) ）。

命题 5. --- 设 \(n \in I\) 。

\begin{enumerate}
\def\labelenumi{\alph{enumi})}
\item
  对每个子空间 \(\mathrm{F} \in \mathcal{F}_{n+1}\) （适应于 \(u\) 的），我们有 \(\lambda_n(u) = r_{\mathrm{F}}(u)\) ；
\item
  对每个子空间 \(F \in F_{n}\) （适应于 u 的），我们有 \(\lambda_{n}(u) = R_{\mathrm{F}}(u)\) 。
\end{enumerate}

设 \(F \in F_{n+1}\) 是适应于 u 的子空间，\((f_i)_{0 \leqslant i \leqslant n}\) 是 F 的一个标准正交基，使得 \(u(f_i) = \lambda_i(u)f_i\) （对所有 i）。对 F 中所有 x，我们有

\[
\langle x \mid u (x) \rangle = \sum_ {0 \leqslant i \leqslant n} \lambda_ {i} (u) | \langle f _ {i} \mid x \rangle | ^ {2} \geqslant \lambda_ {n} (u) \sum_ {0 \leqslant i \leqslant n} | \langle f _ {i} \mid x \rangle | ^ {2} = \lambda_ {n} (u) \| x \| ^ {2},
\]

其中等式当 \(x = f_{n}\) 时成立。这蕴含 \(r_{\mathrm{F}}(u) = \lambda_{n}(u)\) ，从而得到断言 (a)。

设 \(F \in F_{n}\) 是适应于 u 的子空间。闭子空间 \(F^{\circ}\) 非零（因为 \(n = \dim(F) < \dim(E)\) ）且在 u 下稳定；由 u 通过过渡到子空间得到的自同态 \(\widetilde{u}\) （\(F^{\circ}\) 的）是紧且正的。

有 \(\mathrm{Sp}(\widetilde{u})\subset \mathrm{Sp}(u)\) 。我们还有 \(\mathrm{Sp}(\widetilde{u})\subset [0,\lambda_n(u)]\) 。事实上，只需验证 \(\lambda \leqslant \lambda_{n}(u)\) （对所有 \(\lambda \in \mathrm{Sp}_{\mathrm{s}}(\widetilde{u})\) ）。数 \(\lambda\) 这时是 \(u\) 的特征值，故存在 \(j\in I\) 使 \(\lambda = \lambda_j(u)\) 。于是 \(u\) 的相对于 \(\lambda_{j}(u)\) 的特征空间不包含于 F，这蕴含 \(\lambda_{j}(u)\leqslant \lambda_{n}(u)\) 。

假设 \(\lambda_n(u) > 0\) 。\(u\) 的相对于 \(\lambda_n(u)\) 的特征空间于是不包含于 F，从而 \(\lambda_n(u)\) 属于 \(\widetilde{u}\) 的敏感谱。我们推出 \(\sup (\mathrm{Sp}(\widetilde{u})) = \lambda_n(u)\) 。若 \(\lambda_n(u) = 0\) ，由于 \(\widetilde{u}\) 的谱这时化为 \(\{0\}\) ，我们得到同样的结果。

此外，按定义我们有 \(\mathrm{R}_{\mathrm{F}}(u)=\mathrm{R}_{\{0\}}(\widetilde{u})\) ，最后按 I, p. 139, 命题 9, a)，有 \(\mathrm{R}_{\{0\}}(\widetilde{u})=\sup(\mathrm{Sp}(\widetilde{u}))\) 。断言 b) 得证。

命题 6. --- 对所有 \(n \in I\) ，我们有

\[
\lambda_ {n} (u) = \sup _ {\mathrm{F} \in \mathscr {F} _ {n + 1}} r _ {\mathrm{F}} (u) = \inf _ {\mathrm{F} \in \mathscr {F} _ {n}} \mathrm{R} _ {\mathrm{F}} (u).
\]

设 \((e_{n})_{n\in\mathrm{I}}\) 是具有 IV, p. 152 命题 3 的性质的标准正交族。对每个整数 \(n\) （满足 \(1\leqslant n<M+1\) ），设 \(F_{n}\in\mathcal{F}_{n}\) 是由 \((e_{0},\ldots,e_{n-1})\) 生成的 E 的 n 维子空间；按构造，空间 \(F_{n}\) 适应于 u。于是按命题 5，我们有

\[
\lambda_ {n} (u) = r _ {\mathrm{F} _ {n + 1}} (u) = \mathrm{R} _ {\mathrm{F} _ {n}} (u).\tag{2}
\]

设 \(n \in I\) 。设 \(F \in \mathcal{F}_{n+1}\) 。在 \(F\) 上的限制（到 \(F_n\) 上的正交投影的）不是单射，故存在 \(x \neq 0\) （\(F\) 中的，正交于 \(F_n\) 的）。

由于 \(x \in \mathrm{F}_{n}^{\circ}\) ，我们这时有（IV, p. 152, 命题 3）

\[
\begin{array}{l}\langle x\mid u(x)\rangle = \sum_{\substack{m\in \mathrm{I}\\ m\geqslant n}}\lambda_{m}(u)|\langle e_{m}\mid x\rangle |^{2}\\ \leqslant \lambda_{n}(u)\sum_{\substack{m\in \mathrm{I}\\ m\geqslant n}}|\langle e_{m}\mid x\rangle |^{2} = \lambda_{n}(u)\| x\|^{2}. \end{array}
\]

这证明 \(r_{\mathrm{F}}(u)\leqslant\lambda_{n}(u)\) ，从而特别地有不等式

\[
\sup _ {\mathrm{F} \in \mathscr {F} _ {n + 1}} r _ {\mathrm{F}} (u) \leqslant \lambda_ {n} (u).\tag{3}
\]

设 \(F \in F_{n}\) 。在 \(F_{n+1}\) 上的限制（到 F 上的正交投影的）不是单射，故存在向量 \(x \neq 0\) （\(F_{n+1}\) 中的，正交于 F 的）。由于 \(x \in F_{n+1}\) ，我们有（同处）

\[
\begin{array}{c} \langle x \mid u (x) \rangle = \sum_ {0 \leqslant m \leqslant n} \lambda_ {m} (u) | \langle e _ {m} \mid x \rangle | ^ {2} \\ \geqslant \lambda_ {n} (u) \sum_ {0 \leqslant m \leqslant n} | \langle e _ {m} \mid x \rangle | ^ {2} = \lambda_ {n} (u) \| x \| ^ {2} \end{array}
\]

因而 \(\mathrm{R_F}(u) \geqslant \lambda_n(u)\) 。特别地，由此得出

\[
\inf _ {\mathrm{F} \in \mathcal {F} _ {n}} \mathrm{R} _ {\mathrm{F}} (u) \geqslant \lambda_ {n} (u).\tag{4}
\]

鉴于公式 (2)、(3) 与 (4)，命题得证。

